A lei da gravitação de Newton dá o peso a uma distância $r$ do centro da Terra:
\[
W(r)=\frac{GM_E m}{r^2}.
\]

Na superfície ($r=R$), $W_0=mg=10\,\mathrm{N}$, logo $GM_E m=g R^2$. A altitude é $h=r-R$:
\[
W(h)=W_0\left(\frac{R}{R+h}\right)^{2}.
\]

\textbf{Dados:} $R=6{,}371\times10^{6}\,\mathrm{m}$, $W_0=10\,\mathrm{N}$.

\textbf{A $h=10\,\mathrm{m}$:}
\[
\frac{R}{R+h}=\frac{6{,}371\times10^{6}}{6{,}371\times10^{6}+10}=\frac{1}{1+10/(6{,}371\times10^6)}\approx1-1{,}570\times10^{-6}.
\]
\[
W(10)=10\times\left(1-1{,}570\times10^{-6}\right)^{2}\approx10\times\left(1-3{,}14\times10^{-6}\right)\approx9{,}99997\,\mathrm{N}.
\]

A variação relativa: $\Delta W/W_0\approx3{,}1\times10^{-4}\%$. Imperceptível na prática.

\textbf{No Monte Everest ($h=8848\,\mathrm{m}$):}
\[
\frac{R}{R+h}=\frac{6{,}371\times10^{6}}{6{,}379848\times10^{6}}\approx0{,}99861.
\]
\[
W(8848)=10\times(0{,}99861)^{2}=10\times0{,}99723\approx9{,}9723\,\mathrm{N}.
\]

A variação relativa: $\Delta W/W_0\approx0{,}277\%$.

\textbf{Resumo:}
\begin{center}
\begin{tabular}{lll}
\hline
Altitude & Peso & Variação \\
\hline
Superfície & $10{,}0000\,\mathrm{N}$ & --- \\
$h=10\,\mathrm{m}$ & $9{,}99997\,\mathrm{N}$ & $-0{,}0003\%$ \\
$h=8848\,\mathrm{m}$ & $9{,}9723\,\mathrm{N}$ & $-0{,}277\%$ \\
\hline
\end{tabular}
\end{center}